\documentclass[10pt,a4paper]{amsart}
\usepackage[margin=19mm]{geometry}
\usepackage{amsmath,amssymb,mathtools}
\usepackage[scr=rsfs,cal=euler]{mathalfa}
\usepackage{xfrac}
\usepackage[unicode,hidelinks,bookmarks=false]{hyperref}
\newcommand{\ie}{i.e.\ }
\newcommand{\cf}{cf.\ }
\newcommand{\resp}{respectively\ }
\newcommand{\Subsection}{\subsection}
\providecommand{\nobreakdash}{\nobreak\hbox{-}}
\setlength{\emergencystretch}{2em}
\multlinegap0pt
\usepackage{amssymb}
\newtheorem{theorem}{Theorem}
\def\leq{\leqslant}
\title{Determinantal point processes on complex manifolds: Construction and limit theorems}
\author{Thibaut Lemoine}
\date{Source: arXiv:2211.06955v3 (CC BY 4.0)}
\begin{document}
\maketitle

\noindent\emph{Source and licence.} Determinantal point processes on complex manifolds: Construction and limit theorems. Version arXiv:2211.06955v3 (CC BY 4.0), distributed by arXiv under the Creative Commons Attribution 4.0 International licence, http://creativecommons.org/licenses/by/4.0/, as recorded in the arXiv licence metadata for this exact version and shown on the abstract page https://arxiv.org/abs/2211.06955v3. Full licence notice: \texttt{licenses/arxiv-2211.06955-CC-BY-4.0.txt}.\par\smallskip

\noindent\textit{Editorial scope.} Counted unit: the article's theorem thm:intrinsic-determinantal-property (source0.tex lines 575--584). The article's complete original proof of it is delivered with it (source0.tex lines 586--704, one \texttt{proof} environment of 119 lines). No mathematics is written, repaired or re-derived: the statement and the proof are the article's own lines, in the article's own notation, and no retained line is substituted or re-flowed.  No further source-local context is needed: every cross-reference of every retained line resolves inside the two delivered spans.  Nothing was omitted from any retained span, so the package carries no omission record.  No retained line contains a citation, so the wrapper carries no bibliography and no bibliography entry.  No retained line contains a figure environment, an image-inclusion command or drawing code, so no image is delivered and the package declares no asset dependency.  Editorial formatting only, in addition: an AMS \texttt{amsart} preamble that restates the article's own declarations and macros used by the retained lines, namely the environments \texttt{theorem} on separate counters, together with the standard packages those lines need.  The retained environments print in this document as Theorem 1 in order of appearance, whereas the article prints the same items with the numbering of its own full document.  The complete native source of the article as distributed in the arXiv e-print for this exact version is bundled unchanged as \texttt{sources/133415/source0.tex}.
\begin{theorem}\label{thm:intrinsic-determinantal-property}
Let $X$ be the point process associated with $\mathbb P_H$. Then $X$ is determinantal: for every $1\leq m\leq N$, its $m$-point correlation function is
\[
    \rho_m(x_1,\ldots,x_m)
    =
    \det_{\mathbf x}
    \big(B_H(x_i,x_j)\big)_{1\leq i,j\leq m}.
\]
For $m>N$, one has $\rho_m=0$.
\end{theorem}

\begin{proof}
Let $s_1,\ldots,s_N$ be an orthonormal basis of $H$. For every subset $I=\{i_1<\cdots<i_r\}\subset\{1,\ldots,N\}$, define the partial Slater section
\[
S_I(z_1,\ldots,z_r)
=
\sum_{\sigma\in S_r}
\operatorname{sgn}(\sigma)\,
s_{i_{\sigma(1)}}(z_1)\otimes\cdots\otimes s_{i_{\sigma(r)}}(z_r).
\]
Thus $S_{\{1,\ldots,N\}}=S_H$.

We first record two elementary identities. Let $\mathbf x=(x_1,\ldots,x_m)$. By the Cauchy--Binet formula applied to the evaluation map $\operatorname{ev}_{\mathbf x}:H\longrightarrow
L_{x_1}\oplus\cdots\oplus L_{x_m},$ one has
\begin{equation}\label{eq:32}
\det_{\mathbf x}\bigl(B_H(x_i,x_j)\bigr)_{1\leq i,j\leq m}
=
\sum_{\substack{I\subset\{1,\ldots,N\}\\ |I|=m}}
\|S_I(\mathbf x)\|^2.
\end{equation}
Indeed, after choosing unit vectors in the fibers $L_{x_1},\ldots,L_{x_m}$, the matrix of $B_H(\mathbf x)=\operatorname{ev}_{\mathbf x}\operatorname{ev}_{\mathbf x}^*$ is $AA^*$, where $A$ is the $m\times N$ matrix of the evaluation map. Hence
\[
\det(AA^*)
=
\sum_{|I|=m} |\det A_I|^2,
\]
and $|\det A_I|^2$ is precisely $\|S_I(\mathbf x)\|^2$.

Next, for subsets $I,J\subset\{1,\ldots,N\}$ with $|I|=|J|=r$, orthonormality gives
\begin{equation}\label{eq:33}
\int_{M^r}
\left\langle
S_I(\mathbf z),S_J(\mathbf z)
\right\rangle
\,d\mu(z_1)\cdots d\mu(z_r)
=
r!\,\mathbf 1_{\{I=J\}}.
\end{equation}
This follows by expanding both Slater sections and using
\[
\int_M h_x(s_a(x),s_b(x))\,d\mu(x)=\delta_{ab}.
\]

Now fix $1\leq m\leq N$, and write
\[
\mathbf x=(x_1,\ldots,x_m),
\qquad
\mathbf y=(y_1,\ldots,y_{N-m}).
\]
The Laplace expansion of the full Slater section with respect to the first $m$ variables gives
\[
S_H(\mathbf x,\mathbf y)
=
\sum_{\substack{I\subset\{1,\ldots,N\}\\ |I|=m}}
\varepsilon(I)\,
S_I(\mathbf x)\otimes S_{I^c}(\mathbf y),
\]
where $\varepsilon(I)\in\{\pm1\}$ is the sign of the shuffle putting the indices of $I$ in the first $m$ positions and the indices of $I^c$ in the remaining positions.

Taking the squared norm, integrating over $\mathbf y$, and using~\eqref{eq:33}, we obtain
\[
\begin{aligned}
\int_{M^{N-m}}
\|S_H(\mathbf x,\mathbf y)\|^2
\,d\mu(y_1)\cdots d\mu(y_{N-m})
&=
\sum_{I,J}
\varepsilon(I)\varepsilon(J)
\left\langle S_I(\mathbf x),S_J(\mathbf x)\right\rangle
\\
&\qquad\qquad\times
\int_{M^{N-m}}
\left\langle
S_{I^c}(\mathbf y),S_{J^c}(\mathbf y)
\right\rangle
\,d\mu^{\otimes(N-m)}(\mathbf y)
\\
&=
(N-m)!
\sum_{\substack{I\subset\{1,\ldots,N\}\\ |I|=m}}
\|S_I(\mathbf x)\|^2.
\end{aligned}
\]
Combining this with~\eqref{eq:32}, we get
\begin{equation}\label{eq:35}
\frac{1}{(N-m)!}
\int_{M^{N-m}}
\|S_H(\mathbf x,\mathbf y)\|^2
\,d\mu^{\otimes(N-m)}(\mathbf y)
=
\det_{\mathbf x}\bigl(B_H(x_i,x_j)\bigr)_{1\leq i,j\leq m}.
\end{equation}
Since the density of $P_H$ is $\frac{1}{N!}\|S_H(x_1,\ldots,x_N)\|^2$ with respect to $\mu^{\otimes N}$, the $m$-point correlation function is
\[
\begin{aligned}
\rho_m(x_1,\ldots,x_m)
&=
\frac{N!}{(N-m)!}
\int_{M^{N-m}}
\frac{1}{N!}
\|S_H(x_1,\ldots,x_m,y_1,\ldots,y_{N-m})\|^2
\,d\mu(y_1)\cdots d\mu(y_{N-m})
\\
&=
\frac{1}{(N-m)!}
\int_{M^{N-m}}
\|S_H(\mathbf x,\mathbf y)\|^2
\,d\mu^{\otimes(N-m)}(\mathbf y).
\end{aligned}
\]
By~\eqref{eq:35}, this equals
\[
\rho_m(x_1,\ldots,x_m)
=
\det_{\mathbf x}\bigl(B_H(x_i,x_j)\bigr)_{1\leq i,j\leq m}.
\]
This proves the claimed determinantal formula for $1\leq m\leq N$.

Finally, since the process has exactly $N$ particles, its factorial moment measures of order $m>N$ vanish. Hence $\rho_m=0$ for $m>N$.
\end{proof}

\end{document}
